\documentclass[10pt,aspectratio=169]{beamer}
\input{../../_en_preambulo_beamer}% Comandos y macros estadísticas compartidas para las presentaciones Beamer en inglés (Metropolis)

% Conjuntos numéricos básicos
\newcommand{\R}{\mathbb{R}}
\newcommand{\N}{\mathbb{N}}
\newcommand{\Q}{\mathbb{Q}}
\newcommand{\Z}{\mathbb{Z}}
\newcommand{\set}[1]{\left\{ #1 \right\}}
\newcommand{\abs}[1]{\left|#1\right|}
\newcommand{\norm}[1]{\left\| #1 \right\|}

% Comandos estadísticos y probabilísticos del curso
\newcommand{\Var}{\operatorname{Var}}
\newcommand{\cov}{\operatorname{Cov}}
\newcommand{\corr}{\rho}
\newcommand{\comb}[2]{\begin{pmatrix} #1 \\ #2 \end{pmatrix}}
\newcommand{\s}{\sigma}
\newcommand{\particion}{\mathfrak{P}}
\newcommand{\card}[1]{\#\left( #1 \right)}
\newcommand{\seq}[3]{\set{#1}_{#2}^{#3}}
\newcommand{\E}{\mathbb{E}}
\newcommand{\Prob}{\mathbb{P}}

% Entornos pedagógicos y cajas en inglés académico natural (soportando tanto nombres en inglés como en español)
\newenvironment{discussion}{%
  \begin{alertblock}{Discussion Question}%
}{%
  \end{alertblock}%
}
\newenvironment{discusion}{%
  \begin{alertblock}{Discussion Question}%
}{%
  \end{alertblock}%
}

\newenvironment{reflection}{%
  \begin{alertblock}{Classroom Reflection}%
}{%
  \end{alertblock}%
}
\newenvironment{reflexion}{%
  \begin{alertblock}{Classroom Reflection}%
}{%
  \end{alertblock}%
}

\newenvironment{paradox}{%
  \begin{alertblock}{Exploratory Paradox}%
}{%
  \end{alertblock}%
}
\newenvironment{paradoja}{%
  \begin{alertblock}{Exploratory Paradox}%
}{%
  \end{alertblock}%
}

\newenvironment{motivation}{%
  \begin{exampleblock}{Motivation: Data Science in Practice}%
}{%
  \end{exampleblock}%
}
\newenvironment{motivacion}{%
  \begin{exampleblock}{Motivation: Data Science in Practice}%
}{%
  \end{exampleblock}%
}

\newenvironment{quickchallenge}{%
  \begin{block}{Quick Team Challenge}%
}{%
  \end{block}%
}
\newenvironment{retoclase}{%
  \begin{block}{Quick Team Challenge}%
}{%
  \end{block}%
}

\title[Homogeneity tests]{Section 07.11: Homogeneity Tests}\subtitle{One categorical variable across groups}
\author[J. Castillo Colmenares]{Juliho Castillo Colmenares}\institute{Tecnológico de Monterrey}\date{}
\begin{document}
\begin{frame}[plain]\titlepage\end{frame}
\begin{frame}{Roadmap}\small\begin{enumerate}\item Design and concept.\item Frequency table.\item Statistic and decision.\end{enumerate}\end{frame}
\begin{frame}{Source and scope}\small The test compares one categorical distribution across independent subgroups.\begin{block}{Guiding question}Do the groups share the same categorical distribution?\end{block}\end{frame}
\begin{frame}{Homogeneity and independence}\small The chi-square formula is the same as for independence, but the design differs: samples are fixed by population and one variable is compared.\end{frame}
\begin{frame}{Sampling design}\small Draw $c$ independent samples of sizes $n_1,\ldots,n_c$ and observe $r$ categories in each sample.\end{frame}
\begin{frame}{Hypotheses}\small\[H_0:\text{the }c\text{ populations share a distribution};\qquad H_a:\text{at least one differs}.\]\end{frame}
\begin{frame}{Expected frequencies}\small For an $r\times c$ table,\[E_{ij}=\frac{(\text{row total }i)(\text{column total }j)}{\text{grand total}}.\]\end{frame}
\begin{frame}{Chi-square statistic}\small\[\chi^2=\sum_{i=1}^{r}\sum_{j=1}^{c}\frac{(O_{ij}-E_{ij})^2}{E_{ij}}.\]\end{frame}
\begin{frame}{Degrees of freedom}\small Compare with $\chi^2_{(r-1)(c-1)}$. Fixed margins reduce the number of free entries.\end{frame}
\begin{frame}{Conditions}\small Observations must be independent and expected frequencies sufficiently large for the chi-square approximation.\end{frame}
\begin{frame}{Procedure I}\small\begin{enumerate}\item Define groups and categories.\item Build the observed table.\item Compute expected frequencies.\end{enumerate}\end{frame}
\begin{frame}{Procedure II}\small\begin{enumerate}\setcounter{enumi}{3}\item Compute $\chi^2$.\item Obtain a critical value or $p$-value.\item Interpret the distribution, not one cell alone.\end{enumerate}\end{frame}
\begin{frame}{Reading the statistic}\small Cells with $O_{ij}$ far from $E_{ij}$ contribute more to $\chi^2$. The sum aggregates discrepancies across the table.\end{frame}
\begin{frame}{Application: table}\small Arrange frequencies of one categorical variable in rows and independent groups in columns. Column totals reflect the design.\end{frame}
\begin{frame}{Application: expected counts}\small Compute every $E_{ij}$ from marginal totals and check row and column totals.\end{frame}
\begin{frame}{Application: contrast}\small Evaluate $\chi^2$ with $df=(r-1)(c-1)$ at level $\alpha$.\end{frame}
\begin{frame}{Application: decision}\small A small $p$-value gives evidence that not all distributions are equal.\end{frame}
\begin{frame}{What the test does not identify}\small A global rejection does not identify which groups or categories drive the difference. Follow-up comparisons and substantive interpretation are needed.\end{frame}
\begin{frame}{Common errors}\small\begin{itemize}\item Call the homogeneity design independence.\item Ignore small expected counts.\item Treat a global rejection as a one-cell diagnosis.\end{itemize}\end{frame}
\begin{frame}{Reporting template}\small\begin{block}{Report} ``The $r\times c$ table gave $\chi^2=\cdots$, $df=\cdots$, $p=\cdots$. At $\alpha=\cdots$, we [reject/do not reject] homogeneity.''\end{block}\end{frame}
\begin{frame}{Synthesis}\small\begin{itemize}\item Homogeneity compares populations through one categorical variable.\item The statistic uses observed and expected counts.\item A global rejection requires follow-up analysis.\end{itemize}\end{frame}
\begin{frame}{Chapter connection}\small This test is the general framework for the multiple-proportion contrast that follows.\end{frame}
\end{document}
